\documentclass[twocolumn]{aastex631}
\begin{document}
\title{A residual reionization ionizing-photon-budget shortfall for star-forming galaxies at z\~{}9}
\author{NebulaMind Lab (autonomous pipeline)}
\affiliation{NebulaMind Lab, \url{https://lab.nebulamind.net}}
\begin{abstract}
Reionization ionizing-photon-budget reconciliation at z\~{}9: star-forming galaxies require f\_esc=0.390 (+0.393/-0.200) to close the budget (Madau-Dickinson SFRD, log xi\_ion=25.5+/-0.15, clumping C=2-5, JWST-SFRD tail), versus indirect-proxy-inferred f\_esc=0.062 (+0.108/-0.039) from LzLCS O32/beta calibrations. Median delta(required-inferred)=+0.302 dex-frac (16-84\%: +0.087 to +0.697); 93\% of the systematic MC shows a shortfall. Conclusion: a genuine SHORTFALL remains, and the sign holds under both O32 and beta calibrations. The result is bounded by the xi\_ion x clumping x proxy-calibration systematic, not by statistics. Generated autonomously from public data (jwst) via the NebulaMind Lab runner: a bounded, reproducible, descriptive study.
\end{abstract}
\section{Introduction}
Recent studies have highlighted a potential crisis in reconciling the photon budget during reionization [Muñoz2024]. This discrepancy raises questions about our understanding of the ionizing sources responsible for this process. Previous works have emphasized the importance of accurately accounting for the ionizing photons produced by star-forming galaxies and their escape fraction [Davies2021, Park2022]. In light of these findings, it is essential to reassess the reionization photon budget using established literature values.
\section{Data and method}
To address this issue, we adopt a systematic approach grounded in published data. We utilize the Madau \& Dickinson (2014) analytic fitting function for the cosmic star formation rate density (SFRD). For ionizing efficiency (xi\_ion), we rely on previously published values, specifically log xi\_ion = 25.5 ± 0.15 [Madau2017]. Additionally, we employ the O32/beta f\_esc proxy calibrations from LzLCS studies [Chisholm+22, Flury+22; Simmonds+24] to estimate escape fractions.
\section{Result}
Our analysis reveals that star-forming galaxies must have an escape fraction of f\_esc = 0.390 (+0.393/-0.200) at z\~{}9 to reconcile the reionization photon budget. This value is significantly higher than the indirect-proxy-inferred escape fraction of f\_esc = 0.062 (+0.108/-0.039) derived from LzLCS O32/beta calibrations. The median difference between the required and inferred escape fractions is +0.302 dex-frac, with a range of +0.087 to +0.697. Notably, 93\% of our systematic Monte Carlo simulations indicate a shortfall in the photon budget.
\begin{figure}\centering\includegraphics[width=\columnwidth]{result.png}\caption{A residual reionization ionizing-photon-budget shortfall for star-forming galaxies at z\~{}9}\end{figure}
\section{Caveats}
It is crucial to acknowledge the limitations of this study. Our analysis relies on a single selection of literature values and does not incorporate new observational data or account for potential variations in xi\_ion and clumping factor (C) across different environments. Furthermore, the O32/beta proxy calibrations may introduce systematic uncertainties due to their reliance on specific assumptions about galaxy properties. These factors underscore the need for further research and refined measurements to better understand the reionization process and its underlying mechanisms.
\section*{References}
\footnotesize
[Muoz2024] Muñoz et al. (2024). Reionization after JWST: a photon budget crisis?. bibcode 2024MNRAS.535L..37M \par [Davies2021] Davies et al. (2021). The Predicament of Absorption-dominated Reionization: Increased Demands on Ionizing Sources. bibcode 2021ApJ...918L..35D \par [Park2022] Park et al. (2022). Calibrating excursion set reionization models to approximately conserve ionizing photons. bibcode 2022MNRAS.517..192P \par [Duncan2015] Duncan et al. (2015). Powering reionization: assessing the galaxy ionizing photon budget at z < 10. bibcode 2015MNRAS.451.2030D \par [Madau2017] Madau (2017). Cosmic Reionization after Planck and before JWST: An Analytic Approach. bibcode 2017ApJ...851...50M

\end{document}
